\documentclass{article}
\usepackage[T2A]{fontenc}
\usepackage[utf8]{inputenc}
\usepackage[english, russian]{babel}

% Set page size and margins
% Replace `letterpaper' with`a4paper' for UK/EU standard size
\usepackage[a4paper,top=2cm,bottom=2cm,left=2cm,right=2cm,marginparwidth=1.75cm]{geometry}

\usepackage{amsmath}
\usepackage{graphicx}
\usepackage[colorlinks=true, allcolors=blue]{hyperref}
\usepackage{amsfonts}
\usepackage{amssymb}
% \usepackage[left=1cm,right=1cm,top=1cm,bottom=1cm]{geometry}
\usepackage{hyperref}
\usepackage{seqsplit}
\usepackage[dvipsnames]{xcolor}
\usepackage{enumitem}
\usepackage{algorithm}
\usepackage{algpseudocode}
\usepackage{algorithmicx}
\usepackage{mathalfa}
\usepackage{mathrsfs}
\usepackage{dsfont}
\usepackage{caption,subcaption}
\usepackage{wrapfig}
\usepackage[stable]{footmisc}
\usepackage{indentfirst}
\usepackage{rotating}
\usepackage{pdflscape}

\usepackage{MnSymbol,wasysym}
\usepackage{listings}

\lstset{
    language=Python,
    basicstyle=\ttfamily\small,
    keywordstyle=\color{blue},
    stringstyle=\color{red},
    commentstyle=\color{gray},
    backgroundcolor=\color{lightgray!20},
    frame=single,
    rulecolor=\color{black},
    numbers=left,
    numberstyle=\tiny\color{gray},
    stepnumber=1,
    numbersep=10pt,
    showspaces=false,
    showstringspaces=false,
    breaklines=true,
    breakatwhitespace=true,
    tabsize=4,
    captionpos=b
}

\begin{document}

\begin{center}
    \Large {AI Masters \\
        Алгоритмы-2024 \\
        Осень \\
        Семестровая контрольная (midterm) \\
        Никита Загайнов}
\end{center}

\bigskip

\textbf{1[3]} Функции $f_i(n)$ таковы, что $\exists C > 0, N \in \mathbb{N} : \forall n > N \hookrightarrow C_1 n^3 \leq f_i(n) \leq C_2 n^5$

Найдите наилучшие оценки сверху и снизу на функцию $g(n) = \sum\limits_{i=1}^n f_i(n)$

\medskip
\textbf{Решение}

\[
    C_1 n^3 \leq f_i(n) \leq C_2 n^5
\]
\[
    \sum_{i=1}^n C_1 n^3 \leq \sum_{i=1}^n f_i(n) \leq \sum_{i=1}^n C_2 n^5
\]
\[
    C_1 n^3 \sum_{i=1}^n 1 \leq g(n) \leq C_2 n^5 \sum_{i=1}^n 1
\]
\[
    C_1 n^4 \leq g(n) \leq C_2 n^6
\]

Таким образом, наилучшие оценки сверху и снизу для функции $g(n)$:

\[
    C_1 n^4 \leq g(n) \leq C_2 n^6 \Longleftrightarrow g(n) = O(n^6), \ g(n) = \Omega(n^4)
\]

\medskip

\textbf{2[2]} По некоторым современным оценкам общее число живших на Земле людей составляет около 100 миллиардов. Сколько букв должно быть в именах, чтобы никого не перепутать? Считаем, что алфавит на всех один и состоит из 34 букв.

\medskip

\textbf{Решение}

Для существования уникального имени для каждого
человека необходимо, чтобы выполнядись неравенство:

\[
    34^k > 100 \cdot 10^9
\]
Тогда
\[
    k > \log_{34}(100 \cdot 10^9) \approx 7.18
\]
Тогда количество букв в имени должно быть не менее
\[
    k \geq 8
\]
\medskip

\textbf{3[3]} На вход алгоритма поступает $n$-битовое число. Он печатает все нечетные числа от 1 до числа, поступившего на вход. Оцените его время работы, если печать $n$-битового числа занимает $O(n)$.

\medskip

\textbf{Решение}

Рассмотрим следующий алгоритм:
\begin{lstlisting}
def print_odd_nums(n):
    cur_num = 1:
    while cur_num <= n:
        print(cur_num)
        cur_num += 2
\end{lstlisting}
Он выполняет операции сложения и вывода числа
за $O(n)$, и таких операций будет $\frac{n}{2}$,
таким образом, время работы алгоритма
\[
    O(n^2)
\]
\medskip

\textbf{4[4]} Оцените сложность возведения в $n$-битовую степень матрицы $2 \times 2$ с $n$-битовыми элементами по $n$-битовому модулю.

\medskip

\textbf{Решение}

Возведение в степень матрицы $2 \times 2$ с
$n$-битовыми элементами по $n$-битовому модулю
эквивалентно возведению в степень двух собственных
значений матрицы в её спектральном разложении:
\[
    A^k = Q \Lambda^k Q^{-1}
\]
Так как размер матрицы фиксирован, то сложность
нахождения спектрального разложения и перемножения
матриц асимптотически равна самой сложной операции
в этих процессах, а именно умножению чисел длины
$n$ --- $O(n^2)$. Быстрое возведение применяется
к константному количеству чисел, тогда сложность
возведения $n$-битных в $n$-битную степень по $n$-битовому
модулю будет асимптотически равна $O(n^3)$
Тогда общая сложность алгоритма будет равна
\[
    O(n^2 + n^3 + n^2) = O(n^3)
\]

\medskip

\textbf{5[2 + 3]} Для массива $a$ существует такое $i$, что $a[1] < a[2] < \dots < a[i] > a[i+1] > \dots > a[n]$.

\begin{itemize}
    \item Постройте алгоритм, находящий $i$, обоснуйте его корректность и оцените асимптотику. Найдите и обоснуйте нижнюю оценку для этой задачи.
    \item Постройте алгоритм, находящий $j$, для которого достигается максимум $|a[j] - a[j+1]|$, обоснуйте его корректность и оцените асимптотику. Найдите и обоснуйте нижнюю оценку для этой задачи.
\end{itemize}

\medskip

\textbf{Решение}

\begin{itemize}
    \item Эта задача может быть решена алгоритмом
          бинарного поиска:
          Построим алгоритм для нахождения $i$:
          \begin{lstlisting}
def find_i(a):
    r = len(a) - 1
    l = 0
    while l < r:
        m = (l + r) // 2
        if a[m] < a[m + 1]:
            l = m + 1
        else:
            r = m
    return l
          \end{lstlisting}

          На каждом шаге
          алгоритма проверяем, верно ли, что
          $a[mid] < a[mid + 1], \ m = \lfloor \frac{left + right}{2} \rfloor$.
          Если это так, то искомый элемент находится правее
          $mid$, иначе - левее. Таким образом, на каждом шаге
          мы снижаем диапазон поиска вдвое. Таким образом,
          алгоритм работает за
          \[
              O(\log n)
          \]

          Докажем, что невозможно получить
          алгоритм, работающий быстрее, чем за
          $\lceil \log n \rceil$ шагов. Возьмём
          для примера $\lceil \log n \rceil - 1$
          шагов. Тогда на каждом шаге необходимо,
          чтобы алгоритм отбрасывал более, чем половину
          элементов, но это невозможно, так как
          иначе есть риск отбросить меньше половины
          элементов (при делении не ровно пополам).

    \item
          Построим алгоритм для нахождения $j$:
          \begin{lstlisting}
def find_j(a):
    max_diff = 0
    j = 0
    for i in range(len(a) - 1):
        if abs(a[i] - a[i + 1]) > max_diff:
            max_diff = abs(a[i] - a[i + 1])
            j = i
    return j
          \end{lstlisting}

          Его асимптотика равна $O(n)$, так как
          он полностью проходит по массиву.
          Докажем, что это нижняя оценка для этой задачи.
          Эта задача не может быть решена быстрее,
          чем за $O(n)$, так как у нас нет информации
          о дистанции между соседними элементами, и
          нам необходимо рассмотреть все $n-1$ пар
          соседних элементов.

\end{itemize}

\medskip

\textbf{6[1+1+1+1+1]}

\begin{itemize}
    \item $T(n) = T(\frac{n}{2}) + T(\frac{2n}{3}) + n^2$
    \item $T(n) = T(n - 1) + n^2$
    \item $T(n) = 4 T(\frac{n}{16}) + \sqrt{n}$
    \item $T(n) = 9 T(\frac{n}{3}) + n^3$
    \item $T(n) = 9 T(\frac{n}{3}) + n$
\end{itemize}

\medskip

\textbf{Решение}

\begin{itemize}
    \item $T(n) = T(\frac{n}{2}) + T(\frac{2n}{3}) + n^2$ \\

          Очевидно, что
          \[
              T(n) \leq 2T(\frac{2n}{3}) + n^2
          \]
          и \[
              T(n) \geq 2T(\frac{n}{2}) + n^2
          \]
          Для первого соотношения по третьему случаю
          Мастер-теоремы
          получаем верхнюю оценку $\Theta(n^2)$, так как выполняется
          условие
          \[
              f(n) = \Omega(n^{\log_{1.5} 2} + \epsilon)
          \]
          и
          \[
              af(\frac{n}{b}) \leq cf(n) \Longleftrightarrow
              2f(\frac{n^2}{2.25}) \leq \frac{8n^2}{9}
          \]

          Для второго случая также применим третий случай Мастер-теоремы,
          получаем нижнюю оценку $\Theta(n^2)$, так как выполняется
          условие
          \[
              f(n) = \Omega(n^{\log_{2} 2} + \epsilon)
          \]
          и
          \[
              af(\frac{n}{b}) \leq cf(n) \Longleftrightarrow
              2f(\frac{n^2}{4}) \leq \frac{2n^2}{4}
          \]

          Таким образом,
          \[
              T(n) = \Theta(n^2)
          \]

    \item $T(n) = T(n - 1) + n^2$ \\

          \[
              T(n) = T(n - 1) + n^2 =
              \sum_{i=1}^{n} i^2 \geq
              \frac{n}{2} \cdot \frac{n}{2}^2 =
              \frac{n^3}{8} \Theta(n^3) \\
          \]

          \[
              T(n) = \sum_{i=1}^{n} i^2 \leq
              n \cdot n^2 = n^3 = \Theta(n^3)
          \]

          Таким образом, $T(n) = \Theta(n^3)$

    \item $T(n) = 4 T(\frac{n}{16}) + \sqrt{n}$ \\

          К задаче применим второй случай Мастер-теоремы:
          \[
              f(n) = \sqrt{n} = \Theta(n^{\log_{b} n}) = \Theta(\sqrt{n})
          \]
          Таким образом, $T(n) = \Theta(\sqrt{n} \log n)$

    \item $T(n) = 9 T(\frac{n}{3}) + n^3$ \\

          К задаче применим третий случай Мастер-теоремы:
          \[
              f(n) = n^3 = \Omega(n^{\log_{b} a + \epsilon}) =
              \Omega(n^{2 + \epsilon})
          \]
          и
          \[
              af(\frac{n}{b}) \leq cf(n) \Longleftrightarrow
              9f(\frac{n^3}{27}) \leq \frac{n^3}{3}
          \]

          Таким образом, $T(n) = \Theta(n^3)$

    \item $T(n) = 9 T(\frac{n}{3}) + n$ \\

          Применим первый случай Мастер-теоремы:
          \[
              f(n) = n = \Theta(n^{\log_{b} a}) = O(n^2)
          \]
          Таким образом, $T(n) = \Theta(n^2)$
\end{itemize}

\medskip

\textbf{7[3]} Оцените сложность работы алгоритма.

\begin{lstlisting}
def f(n):
    for i = 0; i < n; i += 1:
        for j = 0; j < i; j += 1:
            for m = 1; m < j**2; m *= 2:
                print("hehe")

            for m = 0; m < n; m += 1:
                print("hehe")
\end{lstlisting}

\medskip

\textbf{Решение}

Общее количество итераций равно

\begin{align*}
     & \sum_{i=0}^{n-1} \sum_{j=0}^{i-1}
    \left( \sum_{m=1}^{\lfloor \log j^2 \rfloor} 1 +
    \sum_{m=0}^{n - 1} 1 \right) =                           \\
     & \sum_{i=0}^{n-1} \sum_{j=0}^{i-1}
    \sum_{m=1}^{\lfloor \log{(j+1)^2} \rfloor} 1 +
    \sum_{i=0}^{n-1} \sum_{j=0}^{i-1} \sum_{m=0}^{n - 1} 1 = \\
     & \sum_{i=0}^{n-1} \sum_{j=0}^{i-1}
    \lfloor 2\log(j+1) \rfloor +
    \sum_{i=0}^{n-1} \sum_{j=0}^{i-1} n
\end{align*}

Обозначим последнее выражение за $S_1(n) + S_2(n)$
\[
    S_1(n) \leq \sum_{i=0}^{n-1} (i - 1)
    \lfloor 2\log i \rfloor \leq
    (n-1)(n-1) \lfloor 2\log n \rfloor = \Theta(n^2 \log n)
\]
\[
    S_1(n) \geq \sum_{i=0}^{n-1} \lfloor \frac{i - 1}{2} \rfloor
    \lfloor 2\log \lfloor \frac{i - 1}{2} \rfloor \rfloor \geq
    \lfloor \frac{n-1}{2} \rfloor \lfloor \frac{n-1}{2} \rfloor
    \lfloor 2\log \lfloor \frac{n}{2} \rfloor \rfloor = \Theta(n^2 \log n)
\]

и
\[
    S_2(n) = n \sum_{i=0}^{n-1} \sum_{j=0}^{i-1} 1 =
    n \sum_{i=0}^{n-1} i = n\frac{n(n-1)}{2} = \Theta(n^3)
\]

Тогда
\[
    S_1(n) + S_2(n) = \Theta(n^2 \log n) + \Theta(n^3) = \Theta(n^3)
\]
\medskip

\textbf{8[4]} Есть $n$ монет, среди них одна фальшивая, она легче остальных, и чашечные весы. С помощью весов за один ход можно взвешивать две любые группы монет. Они могут давать три ответа: левая чашечка тяжелее, правая, равны. Постройте алгоритм, находящий фальшивую монету за асимптотически оптимальное число шагов, и докажите нижнюю оценку на эту задачу.

\medskip

\textbf{Решение}

Рассмотрим алгоритм, который делит кучу монет
на две равные части и взвесим их с помощью весов.
Отбрасываем ту часть, которая тяжелее, и проводим
это же действие с оставшимися монетами. Повторяем, пока
не останется одна монета. Таким образом, алгоритм имеет
асимптотику
\[
    T(n) = T(\frac{n}{2}) + O(1) = O(\log n)
\]

Докажем, что это нижняя оценка для этой задачи.
Каждую монету можно представить в виде битовой строки
длины $\lceil \log n \rceil$, на шаге $i$ алгоритм
за константное время может может получить один бит
информации, поделив оставшуюся кучу монет на две равные
части. Таким образом, алгоритм должен сделать минимум
$\lceil \log n \rceil$ шагов, чтобы найти фальшивую монету.

\medskip

\textbf{9[3]} Экономный математик варит первую воду на киселе, вторую, ..., седьмую воду на киселе, ..., $n$-ю воду на киселе, и наливает каждую из них в банку. В каждой следующей воде на киселе концентрация киселя (действительное число) строго меньше, чем в предыдущей, и она написана на банке. Некто переставил все банки местами, и теперь они не упорядочены. Постройте алгоритм, позволяющий найти среднюю концентрацию киселя в банках с $k$-й воды на киселе по $m$-ю воду на киселе ($k < m$, нужно найти среднее по $k$-й, $k+1$-й, ..., $m$-й).

\medskip

\textbf{Решение}

Построим следующий алгоритм, беря во внимание, что функция
$median\_of\_medians$ за линейное время находит $k$-ую
порядковую статистику:

\begin{lstlisting}
def find_mean(k, m, a):
    kth = median_of_medians(a, k)
    mth = median_of_medians(a, m)
    
    running_sum = 0
    count = 0
    for elem in a:
    if elem <= kth and elem >= mth:
        running_sum += elem
        count += 1

    return running_sum / count
\end{lstlisting}

Этот алгоритм работает за $O(n)$, учитывая, что
$median\_of\_medians$ работает за $O(n)$ и он один раз
проходится по массиву.

Докажем, что асимптотика этой задачи не может быть быстрее.
Примем во внимание тот факт, что нам необходимо найти
среднее $m - k + 1$ элементов. Это выполнимо только за
$O(n)$, так как нам пройтись по всем слагаемым.

\end{document}